% =====================================================================
%  MODELLO: DOCUMENTO  (Norme di Progetto, Piano di Progetto, Piano di
%  Qualifica, Analisi dei Requisiti, Glossario, Specifica Tecnica,
%  Manuale Utente, documenti della candidatura...)
%
%  1. Create una cartella per il documento e copiateci questo file:
%       <Milestone>/documenti/interni/<nome>/<nome>.tex   (uso interno)
%       <Milestone>/documenti/esterni/<nome>/<nome>.tex   (uso esterno)
%     es.  RTB/documenti/interni/norme-di-progetto/norme-di-progetto.tex
%     (Milestone = Candidatura, RTB oppure PB). Dopo il merge la CI
%     mette accanto il PDF con lo stesso nome; sul sito il nome ha in
%     fondo la versione. \uso qui sotto deve dire come la cartella.
%  2. Cambiate i campi qui sotto (dopo il simbolo % c'è un commento).
%  3. Scrivete il contenuto dopo \begin{document}.
%  Copertina, registro, indice, elenco di figure e tabelle (se ce ne
%  sono), intestazioni e numeri di pagina si fanno da soli.
%  Guida completa: GUIDA.md (ricetta 4).
% =====================================================================
\documentclass[documento]{gpt}

\titolo{Ciao}        % es. Norme di Progetto
\uso{interno}                      % "interno" (es. Norme di Progetto)
                                   % oppure "esterno" (tutti gli altri)
\responsabile{spagnoli}            % chi approva il documento: il Responsabile

% Facoltativi (togliete il % davanti se servono):
% \capitolato{Cn -- Nome del capitolato (Azienda)}   % solo dopo l'assegnazione
% \destinatari{Prof.~Tullio Vardanega\newline Prof.~Riccardo Cardin\newline Azienda}

% ---------------------------------------------------------------------
%  REGISTRO DELLE MODIFICHE  --  la riga più recente va IN CIMA.
%  La versione del documento è quella della PRIMA riga. Redattori e
%  verificatori in copertina si ricavano da qui, da soli.
%
%  \modifica{versione}{data}{autore}{verificatore}{cosa è cambiato}
%
%  Esempio di registro "cresciuto" (la riga nuova si aggiunge sopra):
%  \modifica{1.1.0}{2026-11-20}{mazzurana}{corro}{Aggiunta sezione sui test}
%  \modifica{1.0.0}{2026-11-10}{guglielmo}{corro}{Verifica e approvazione}
%  \modifica{0.2.0}{2026-11-05}{guglielmo}{--}{Stesura della sezione 2}
% ---------------------------------------------------------------------
\modifica{0.1.0}{2026-10-26}{guglielmo}{--}{Prima stesura: struttura del documento}

\begin{document}

\section{Introduzione}

\subsection{Scopo del documento}
A cosa serve questo documento e a chi si rivolge.

\subsection{Glossario}
I termini con la \glo{G} a pedice sono spiegati nel Glossario.

\section{Primo argomento}
Testo.

% Documento lungo? Mettete ogni capitolo in un file a parte nella
% sottocartella capitoli/ e richiamatelo così (senza ".tex"):
% \input{capitoli/secondo-argomento}

\end{document}
